\section{Marco Teórico}

\subsection{Antecedentes externos}

La configuración interleaved en convertidores DC-DC bidireccionales se ha consolidado como una técnica efectiva para mejorar el rendimiento energético y reducir las pérdidas por conmutación. Al operar varias fases desfasadas, esta topología permite reducir la frecuencia de conmutación percibida en cada inductor, lo que disminuye el estrés térmico y eléctrico sobre los componentes. Investigadores en China, como se presenta en \cite{liu2023interleaved}, demostraron que un convertidor buck interleaved de tres fases con inductores acoplados operando a 500 kHz puede alcanzar una eficiencia superior al 98\%, gracias a la reducción del rizado de corriente y la mejora en la distribución de las cargas entre fases. De manera similar, en India, Banerjee et al.~\cite{banerjee2024sic} estudiaron un convertidor buck interleaved de dos fases que utiliza MOSFETs de carburo de silicio (SiC) junto con inductores acoplados, evidenciando una notable reducción en el rizado de corriente y una mejor respuesta dinámica del sistema.\\
El impacto de la configuración interleaved también se observa en la reducción del valor de inductancia requerida y en la mejora del control del rizado de corriente. Nuevamente en China, Liu et al.~\cite{liu2023interleaved} desarrollaron un modelo que relaciona el coeficiente de acoplamiento en inductores integrados con la inductancia equivalente, mostrando cómo un mayor acoplamiento magnético contribuye a minimizar el rizado y a mejorar la eficiencia global del convertidor. De manera similar, en Estados Unidos, en ~\cite{duan2023multiphase} analizaron un convertidor DC-DC bidireccional interleaved con fases variables, donde el ajuste del número de fases y el ciclo de trabajo permite optimizar el rizado de corriente, mejorando la calidad del flujo energético y disminuyendo pérdidas. Estos estudios internacionales respaldan la aplicación de configuraciones interleaved para lograr convertidores más eficientes, compactos y con mejor rendimiento térmico y eléctrico.


\subsection{Antecedentes locales}

En Colombia, se han realizado aportes significativos en el control de convertidores DC-DC bidireccionales. Por ejemplo, en \cite{narvaez2024hess} se presenta un tesis, la cual se destaca por su enfoque en el diseño y validación de algoritmos de control digital adaptativo para convertidores bidireccionales de puente medio, enfatizando la robustez ante variaciones en las condiciones de carga y las perturbaciones externas. Su trabajo incluye la implementación de lazos de control basados en técnicas avanzadas de control predictivo y adaptativo, orientados a optimizar la respuesta dinámica y la estabilidad del sistema bajo escenarios reales de operación.\\
Por otro lado, en \cite{yate2024control} se profundiza en el desarrollo de estrategias de control digital para convertidores DC-DC bidireccionales en configuración interleaved, utilizando plataformas Hardware-in-the-Loop (HIL) para validar el desempeño de los controladores propuestos. Esta investigación pone especial énfasis en la coordinación de fases y la minimización del rizado de corriente mediante la modulación del ciclo de trabajo y el control en tiempo real, aspectos fundamentales para extender la vida útil de los dispositivos de potencia y mejorar la eficiencia energética en sistemas de almacenamiento de energía.\\
Finalmente, una contribución adicional proviene de \cite{gomez2023control} de la Universidad Nacional de Colombia, quienes desarrollaron un modelo de control híbrido para convertidores bidireccionales que combina control clásico y técnicas de inteligencia artificial para la optimización adaptativa del ciclo de trabajo en convertidores interleaved. Su estudio incluye análisis de estabilidad y validación experimental en prototipos físicos, lo que aporta un marco práctico para la aplicación de control avanzado en entornos reales de microrredes DC.
